% Source: https://artofproblemsolving.com/community/c3h1787424_geoaalgebra
% Stable identity: aops-topic-1787424-q01
A circle of area $\pi r^2$ has radius $r$, so the two circles have radii $4$ (since $\pi\cdot4^2=16\pi$) and $10$ (since $\pi\cdot10^2=100\pi$). Their average is $\dfrac{4+10}{2}=7$, and the circle of radius $7$ has area $\pi\cdot7^2=\boxed{49\pi}$.

\begin{asy}
usepackage("amsmath");
size(5cm);
defaultpen(fontsize(10pt));
pair O=(0,0);
draw(circle(O,10),linewidth(0.8));
draw(circle(O,7),blue+linewidth(0.9));
draw(circle(O,4),linewidth(0.8));
dot(O);
draw(O--4*dir(200),linewidth(0.6));draw(O--7*dir(90),blue+linewidth(0.6));draw(O--10*dir(-30),linewidth(0.6));
label("$4$",2*dir(200),dir(110));
label("$7$",(0,4.6),E,blue);
label("$10$",8.5*dir(-30),dir(60));
label("$O$",O,dir(-115));
label("\mbox{}",(min(currentpicture,user=true)+max(currentpicture,user=true))/2,black); // keeps the text after the figure black when LaTeX typesets the labels inline
\end{asy}

(Note that this is not the average of the two areas, $58\pi$.)
